\section{Conclusions and Future Work}
\label{sec:conclusions}

\subsection{Summary}

This project designed, built, documented, and evaluated a Security
Operations Centre home lab using entirely free and open-source tooling.
The lab consists of a Wazuh SIEM manager and two monitored endpoints ---
one Windows and one Linux --- on an isolated virtualised network. Both
endpoints run deliberate detection-grade telemetry: Sysmon with the
SwiftOnSecurity baseline on Windows, and the Linux Audit framework with
\texttt{execve} monitoring on Linux. Four attack simulations covering
persistence, privilege escalation, credential access, and detection
configuration were executed and triaged against the built environment,
each producing either a documented true-positive alert cluster or a
deliberately documented negative result. Full analyst-style incident
reports were written for every simulation, and the entire build was
recorded as a reproducible step-by-step guide alongside a formal report
(this document).

The seven project objectives set out in Section~\ref{sec:aim} were all
met, and the four success criteria stated in the Initial Project Overview
were all satisfied.

\subsection{Answers to the Research Questions}

\textbf{RQ1.} A credible SOC monitoring environment can be built with only
free and open-source tools, provided that "credible" is scoped to
individual learning and portfolio demonstration rather than to running a
production security programme. The Wazuh stack in particular is genuinely
representative of a live SIEM in both operation and concept.

\textbf{RQ2.} Meaningful detection requires deliberate configuration
beyond defaults on both endpoint platforms: Sysmon plus a Wazuh-side
\texttt{<localfile>} block for Windows, and \texttt{auditd} with
\texttt{execve} rules for Linux. In addition, real-time File Integrity
Monitoring on security-critical paths requires an explicit
\texttt{realtime="yes"} directive, without which FIM operates on a
twelve-hour scan cycle. Vendor defaults optimise for coverage with low
overhead, not for security-critical corner cases, and closing those
corner cases is a large share of what detection engineering is.

\textbf{RQ3.} A single-machine virtualised lab represents perhaps twenty
per cent of a professional SOC's operational surface, and does so along
the axis most relevant to the Tier~1 analyst role: the first sixty
seconds after an alert appears. Scale, team workflow, network detection,
cloud telemetry, threat intelligence enrichment, and response automation
are all absent, and this is documented openly in the Discussion rather
than glossed over.

\subsection{Contribution}

The specific contribution of this project is a fully documented,
reproducible SOC home lab build using open-source tools, paired with a
set of triaged attack simulations that include both positive and
deliberately negative results, and with analyst-style incident reports
for each simulation. The project is deliberately positioned as the first
stage of a two-part progression: it establishes the monitoring foundation
on which a planned second project, focused on detection engineering with
MITRE ATT\&CK and Atomic Red Team, will build.

\subsection{Future Work}

Several extensions were identified during the project and are noted here
for future development. They fall into three groups.

\textbf{Detection engineering (Project~2).} The planned follow-on project
will focus on custom rule development, mapping detection coverage against
MITRE ATT\&CK, and using Atomic Red Team to execute a broader and more
systematic set of attack techniques. The two specific detection gaps
identified in this project --- the absence of an aggregate rule for
repeated failed authentication (Simulation~3) and the lack of real-time
FIM on the Windows hosts file (Simulation~4) --- are both natural starting
points for that work.

\textbf{Additional telemetry sources.} The current lab is host-based only.
Adding network-tier detection through Suricata or Zeek integrated with
Wazuh would extend coverage to lateral movement and command-and-control
traffic that host telemetry alone cannot see. Enabling the Wazuh
vulnerability-detector on a schedule (with the disk consumption issue
discussed in Section~\ref{sec:discussion} mitigated) would add
vulnerability visibility to the existing detection surface. Integrating a
threat intelligence feed for IOC enrichment would move the lab closer to
the enrichment workflows real SOCs depend on.

\textbf{Response automation and case management.} All response actions in
the current lab were performed by hand. Adding Wazuh's active-response
module, or integrating a lightweight SOAR platform such as Shuffle for
playbook automation, would demonstrate the response side of the SOC
lifecycle in the same practical way that this project demonstrates the
detection side. Adding a lightweight case management or ticketing layer
would let the analyst workflow be exercised end to end.

\subsection{Closing Note}

The core aim of this project was to close, for the author personally, the
gap between theoretical knowledge of SIEM operations and practical
hands-on experience of it. That gap has been closed, and the resulting
build and documentation are offered as both a self-contained portfolio
artefact and as a foundation for further work in detection engineering.
The environment stands ready for that follow-on work; the disciplines it
required to build --- deliberate telemetry design, structured triage,
honest documentation, and the willingness to run and record negative
tests --- are the same disciplines the follow-on work will demand at
greater depth.
